\documentclass{article}

\usepackage{circuitikz}
\usepackage[T1]{fontenc} 
\usepackage[UTF8]{inputenc}
\usepackage{amsmath}
\usepackage{amssymb}
\usepackage{fancyhdr}
\usepackage{graphicx}
\usepackage{hyperref}
\usepackage{tikz}
  \usetikzlibrary{arrows}
  \usetikzlibrary{shapes}
  \usetikzlibrary{arrows.meta,topaths}
  \usetikzlibrary{bending}
  \usetikzlibrary{calc}
\usepackage{anyfontsize}
\usepackage{sectsty}
\usepackage{../assets/scripts/tex/color-env}
\usepackage{anyfontsize}
\usepackage{xcolor}
\definecolor{DarkGreenBlue}{HTML}{264653}
\definecolor{LightGreenBlue}{HTML}{2A9D8F}
\definecolor{LightOrange}{HTML}{E9C46A}
\definecolor{DarkOrange}{HTML}{F4A261}
\definecolor{RedOrange}{HTML}{E76F51}
\definecolor{BrightRed}{HTML}{D62828}
\definecolor{DeepBlue}{HTML}{003049}



\usepackage[ngerman]{babel}
\title{Elektrotechnik 1 - Praktikum 3}


\usepackage[
  includehead,
  headheight = 17mm,
  footskip = \dimexpr\headsep+\ht\strutbox\relax,
  tmargin = 0mm,
  bmargin = \dimexpr17mm+2\ht\strutbox\relax,
]{geometry}





\pagestyle{fancy}
\fancyhead[L]{\leftmark}
\fancyhead[R]{}
\fancyfoot[L]{}
\fancyfoot[C]{\thepage}
\fancyfoot[R]{\includegraphics[scale=0.2]{../assets/images/haw.jpg}}
\renewcommand\headrulewidth{0.5pt}


\begin{document}


\thispagestyle{empty}
\begin{tikzpicture}[remember picture,overlay]

  \fill[DeepBlue] (current page.south west) rectangle (current page.north east);

  \begin{scope}

    \foreach \i in {2.5,...,22}
      {
        \node[rounded corners, DeepBlue!90,draw ,regular polygon, regular polygon sides=6, minimum size=\i cm, ultra thick] at ($(current page.west)+(2.5,-5)$) {} ;
      }

  \end{scope}

  \node[rounded corners,fill=DeepBlue!95,text =DeepBlue!5,regular polygon,regular polygon sides=6, minimum size=2.5 cm,inner sep=0,ultra thick] at ($(current page.west)+(2.5,-5)$) {\LARGE \bfseries 2020};

  \foreach \i in {0.5,...,22}
    {
      \node[rounded corners,DeepBlue!90,draw,regular polygon,regular polygon sides=6, minimum size=\i cm,ultra thick] at ($(current page.north west)+(2.5,0)$) {} ;
    }

  \foreach \i in {0.5,...,22}
    {
      \node[rounded corners,DeepBlue!98,draw,regular polygon,regular polygon sides=6, minimum size=\i cm,ultra thick] at ($(current page.north east)+(0,-9.5)$) {} ;
    }

  \foreach \i in {12}
    {
      \node[fill = DeepBlue,rounded corners,draw=DeepBlue,regular polygon,regular polygon sides=6, minimum size=\i cm,ultra thick] at ($(current page.south east)+(-0.2,-0.45)$) {} ;
    }


  \foreach \i in {21,...,6}
    {
      \node[DeepBlue!95,rounded corners,draw,regular polygon,regular polygon sides=6, minimum size=\i cm,ultra thick] at ($(current page.south east)+(-0.2,-0.45)$) {} ;
    }

  \node[left,DeepBlue!5,minimum width=0.625*\paperwidth,minimum height=3cm, rounded corners] at ($(current page.north east)+(0,-9.5)$){{\fontsize{25}{30} \selectfont \bfseries ET2 - Praktikum 5}};

  \node[left,DeepBlue!10,minimum width=0.625*\paperwidth,minimum height=2cm, rounded corners] at ($(current page.north east)+(0,-11)$){{\huge \textit{Resonanz}}};

  \node[left,DeepBlue!5,minimum width=0.625*\paperwidth,minimum height=2cm, rounded corners] at ($(current page.north east)+(0,-13)$){{\Large \textsc{Florian Tietjen\hspace{0.5cm}Eric Antosch}}};

\end{tikzpicture}

\newpage
\thispagestyle{empty}

\tableofcontents


\newpage


\section{Vorbereitung}
Als Vorbereitung der folgenden Versuche werden vorher einige Kenngrößen eines Serienresonanzkreies berechnet.
Die Serienschaltung besteht dabei aus einer Induktivität L = 100 mH und einem Gleichstrom-Drahtwiderstand $R_L = 10 \Omega$ . Der Sinus-Generator
hat einen Innenwiderstand von $R_i = 50\Omega$. Die angegebene Größe der Induktivität gilt nur für Frequenzen f < 1kHz.

\subsection{Berechnung der erforderlichen Kapazität}
Die Formel zur Berechnung der Resonanzfrequenz wird nach der Kapazität C umgestellt.
\begin{align*}
  \omega_r &= \frac{1}{\sqrt{LC}}\\
  C &= \frac{1}{{w_r}^2L}
\end{align*}
Mit der Formel werden nun die Kapazitäten verschiedener Frequenzen ausgerechnet:
\begin{center}

  \begin{tabular}{|c|c|c|c|c|c|}    
    \hline
    Freqeunz               & $100 Hz$                   & $500Hz$                   & $1kHz$            & $5kHz$          & $10kHz$              \\
    \hline
    Kapazität              & $25,33\mu F$                      & $1,013\mu F$                     & $253,3nF$            & $10,13nF$           & $2,53nF$              \\
    \hline
  \end{tabular}

\end{center}

\subsection{Berechnung der Gütefaktoren und Bandbreiten}
Nun werden die zugehörigen Gütefaktoren Q und Bandbreiten $\Delta$f für jeweils den Gesamtverlustwiderstand R$_i$ + R$_L$ und den Spulen-Drahtwiderstand R$_L$ allein berechnet.
Für den Gesamtverlustwiderstand gilt:
\begin{align*}
  R_{ges} &= R_i + R_L= 50\Omega + 10\Omega = 60\Omega
\end{align*}
Die Formel der Güte einer Serienschaltung lautet:
\begin{align*}
  Q &= \frac{1}{R}\cdot\sqrt{\frac{L}{C}}
\end{align*}
Für die Bandbreite $\Delta$f ermittelt sich aus:
\begin{align*}
  B &= \frac{f_r}{Q}
\end{align*}
Für verschiedene Frequenzen ergeben sich dann folgende Güten und Bandbreiten:

\begin{center}

  \begin{tabular}{|c|c|c|c|c|c|}    
    \hline
    Freqeunz               & $100Hz$                   & $500Hz$                   & $1kHz$            & $5kHz$          & $10kHz$              \\
    \hline
    $Q_{R_{ges}}$              & $1,05$               & $5,24$                     & $10,47$            & $52,37$           & $104,78$              \\
    \hline
    $Q_L$              & $6,28$                   & $31,42$                   & $62,83$            & $316,02$          & $628,7$              \\
    \hline
    $\Delta f_{R_{ges}}$              & $95,24$                      & $95,42$                     & $95,51$            & $95,47$           & $95,44$              \\
    \hline
    $\Delta f_L$              & $15,92$                      & $15,91$                     & $15,92$            & $15,82$           & $15,91$              \\
    \hline
  \end{tabular}

\end{center}



\newpage
\section{Resonanz - Eingeschwungener Zustand}
\begin{task}
  IIm ersten Versuch soll der Amplitudengang von $\hat{u}_C$ in der abgebildeten Schaltung mit dem Oszilloskop ermittelt werden. Dazu soll beachtet werden,
  dass $R_i = 50\Omega$ und $\hat{u}_{gen} = 1V$ sind. Die Kapazitätsdekade wird mit den berechneten Werten aus der Vorbereitung eingestellt.
\end{task}
\begin{figure}[h]
  \begin{center}

    \caption{Schaltplan zur Bestimmung der Serienimpedanz eines Drahtwiderstands}
  \end{center}
\end{figure}
\begin{devlist}
  T
  \begin{itemize}
    \item Oszilloskop: Tektronix MDO3012
    \item Funktionsgenerator: Keysight 33210A
    \item HP 4294A Impedance Analyzer
    \item Multimeter: MetraHit X-TRA Multimeter
  \end{itemize}
\end{devlist}

\subsection{Bestimmung von $\hat{u}_C$/$\hat{u}_{gen}$}

\begin{table}[h]
  \begin{center}

    \begin{tabular}{|c|c|c|}
      \hline
          & $500Hz; C=1,01\mu F$ & $1kHz; C= 253,3 nF$ \\
      \hline
      -95 &                      &                     \\
      \hline
      -95 &                      &                     \\
      \hline
      -95 &                      &                     \\
      \hline
      -95 &                      &                     \\
      \hline
      -95 &                      &                     \\
      \hline
      -95 &                      &                     \\
      \hline
      -95 &                      &                     \\
      \hline
      -95 &                      &                     \\
      \hline
      -95 &                      &                     \\
      \hline
      -95 &                      &                     \\
      \hline
      -95 &                      &                     \\
      \hline
      -95 &                      &                     \\
      \hline
      -95 &                      &                     \\
      \hline
    \end{tabular}
    \caption{Messwerte für Versuch 1.1}
    \label{tab:MV}
  \end{center}
\end{table}




\subsection{Auswertung}

Unsere abgelesene Resonanzfrequenz liegt nun also bei: $f_{1Res} = 835kHz$.
Aus dieser können wir mit (\ref{eq:LS}) die gesuchte Serieninduktivität des Drahtwiderstands berechnen: $L_S = 3,633\mu H$\\[3pt]
Um nun die Serienresonanzfrequenz mit der Parallelresonanzfrequenz zu vergleichen, setzen wir unsere errechnete Induktivität in
(\ref{eq:f2}) ein und erhalten dadurch: $f_{2Res} = 648,745kHz$. Wir testen nun mit dem Oszilloskop diese berechnete Resonanzfrequenz, indem wir
die beiden Signale wieder in Phase bringen, wobei wir $f_{2Res} = 652kHz$ ablesen.

Da unsere Rechnungen auf der Annahme beruhen, dass unsere Bauteile, im Besonderen der Kondensator, bei jeder Frequenz die gleichen Kenngrößen haben,
erkennen wir eine mögliche Fehlerquelle. Mit dem HP4294A messen wir bei den bestimmten Resonanzfrequenzen die Kapazität aus und erhalten bei $f_{1Res}$ eine Kapazität
von $C_{f1}=9,53nF$ und bei $f_{2Res}$ eine Kapazität von $C_{f2} = 9,55nF$. Aus den angepassten Werten ergibt sich dann eine Serieninduktivität von $L_S = 3,812\mu F$
und eine Parallelresonanzfrequenz von $f_{2Res} = 666,922kHZ$. Da allerdings diese Messwerte sich implizit auf den "falschen" Messwerten beruhen, sind hier die Toleranzen aufgrund der
Fehlerfortpflanzung größer. Dazu kommen noch Ungenauigkeiten bei den Geräten und deren Bedienung.


\newpage
\section{Resonanz - Schaltverhalten}
\begin{task}
  IIn diesem Versuch wird wieder die Schaltung aus 1.1 aufgebaut. Für die Messungen wird zwischen dem Generator und der Spule/C-Dekade eine Präzisions-Widerstandsdekade verschaltet.
  Die Kreisresonanz wird mit der C-Dekade auf f = 1kHz eingestellt. An die Schaltung wird eine symmetrische Rechteckspannung mit einer Frequenz f = 50Hz, Tastverhältnis a = 0,5 und einem Scheitelwert von $\hat{u}_{gen} = 1,0 V$ angelegt.
  Der zeitliche Verlauf der Spannung am Kondensator wird mittels Oszilloskop dargestellt
\end{task}

\begin{devlist}
  T
  \begin{itemize}
    \item Oszilloskop: Tektronix MDO3012
    \item Kapazitätsdekade Time Electronics Model 1071 CapBox
    \item Funktionsgenerator: Keysight 33210A
    \item Multimeter: MetraHit X-TRA Multimeter
  \end{itemize}
\end{devlist}

\subsection{Serienschwingkreis mit Dämpfung}

Abbildung gedämpfter Schwingkreis\\\\

Messwerte:
\begin{table}[h]
  \begin{center}

    \begin{tabular}{|c|c|c|c|}
      \hline
        $T_d$ &$\overline{u}$ & $y_i$  & $y_{i+1}$\\
      \hline
             ca 1ms&    ca -1v&   ca860mV  &    ca340mV    \\         
      \hline
    \end{tabular}
  \end{center}
\end{table}
Berechnung des Dämpfungsmaß:
\begin{align*}
  \delta &= \frac{1}{T_d} \cdot \ln \left(\frac{Y_i}{Y_{i+1}}\right)\\
  \omega &= 2\pi \cdot f = 2\pi \cdot \frac{1}{T_d} = hier muss ca das gleiche wie unten rauskommen \\
  \omega_r &= \frac{1}{\sqrt{LC}} = \frac{1}{\sqrt{100mH \cdot 253,3nF}} = 6283 s^{-1}
\end{align*}


\subsection{Aperiodischer Grenzfall}

Abbildung Aperiodischer Grenzfall

Wert liegt circa bei 1,0753k Ohm

\end{document}